% Source observation: analyses/018-2026-09-08-results-materials/observations/O007-kernel-realization.md
% Replacement prose for kernel-autoresearch.tex; Analysis 018 O007 and figure_data.json:runtime.
% Source reduction: evidence.py:runtime_subset; figure: plots.py:kernels.
% Underlying implementation audit: Run 029 observations/06-implementation-and-claim-audit.md.
\subsection{Realizing sparsity with specialized inference kernels}
\label{sec:kernel-autoresearch}

Logical sparsity helps execution only when skipping work costs less than
performing it. We evaluate 42 proposals from a human-guided Codex search
against native PyTorch scaled dot-product attention (SDPA) in CUDA graphs
on one RTX5090. The workload is BF16, batch-one, uncached 2,048-token
inference with full vocabulary logits. A proposal qualifies only if its
logits and loss agree with the native reference within specified tolerances
on complete validation (Appendix~\ref{app:kernel-results}).

The best qualified incumbent reaches $1.7830\times$ on the fixed A7-OL1,
$\kappa=0.5$ search checkpoint (Figure~\ref{fig:kernel-autoresearch}, left).
The selected implementation K050 qualifies on all 30 checkpoints in the
14M cohort, with geometric mean speedup $1.2340\times$. Speedup is positively
associated with sparsity (right; OLS $R^2=0.8167$), but checkpoint quality,
topology and weights vary; this is not an equal-quality acceleration comparison.

Relative to the same fusion with sparse paths disabled, K050 is
$1.311\times$ as fast at the search checkpoint and $1.043\times$ as fast
geometrically across the cohort, helping 14 of 30 checkpoints.
Attention skipping is unprofitable:
disabling it improves every checkpoint, raising the mean to $1.2506\times$.
Appendix~\ref{app:kernel-results} details the implementation and ablations.
The full-model gain therefore combines fusion with a conditional benefit
from sparse paths; a large zero-product fraction does not make every
sparse operation worthwhile.

\begin{figure}[!htbp]
\centering
\includegraphics[width=\linewidth]{figures/07-kernel-realization.pdf}
\caption{\textbf{Sparse kernel auto-research for full-model acceleration.}
(a) Best qualified incumbent on fixed A7-OL1, $\kappa=0.5$, across 42
iterations, starting at native $1\times$ and reaching $1.7830\times$.
(b) The selected K050 implementation on 30 Pythia-14M checkpoints; symbols identify recipes. The dashed
line is a descriptive OLS fit with intercept ($R^2=0.8167$); dotted lines
mark native $1\times$. Speedup uses each checkpoint's native SDPA CUDA-graph
baseline on one RTX5090: BF16, batch one, uncached 2,048-token inference and
full vocabulary logits. Qualification covers all 338 validation blocks. All checkpoints
qualify, with geometric mean $1.2340\times$. Here ``best'' refers to the
searched proposals; the attention-dense ablation is faster on every checkpoint
(Table~\ref{tab:kernel-ablations}). Sparsity uses FP16 measurements, while
timing uses BF16; both speedup axes have restricted ranges. The fit describes
association across different checkpoints, not a causal sparsity effect.}
\label{fig:kernel-autoresearch}
\end{figure}
\FloatBarrier
